\documentclass[11pt,a4paper]{article}

\usepackage[margin=2.3cm]{geometry}
\usepackage[T1]{fontenc}
\usepackage{lmodern}
\usepackage[utf8]{inputenc}
\usepackage{microtype}
\usepackage{xcolor}
\usepackage{enumitem}
\usepackage{booktabs}
\usepackage{longtable}
\usepackage{array}
\usepackage{titlesec}
\usepackage{fancyhdr}
\usepackage[colorlinks=true,linkcolor=accent,urlcolor=accent,citecolor=accent]{hyperref}

\definecolor{accent}{RGB}{20,70,120}
\definecolor{soft}{RGB}{240,244,248}

\setlist[itemize]{leftmargin=1.3em,itemsep=2pt,topsep=3pt,parsep=0pt}
\setlist[enumerate]{leftmargin=1.5em,itemsep=2pt,topsep=3pt,parsep=0pt}
\setlength{\parindent}{0pt}
\setlength{\parskip}{5pt}
\renewcommand{\arraystretch}{1.25}

\titleformat{\section}{\Large\bfseries\color{accent}}{\thesection}{0.6em}{}
\titleformat{\subsection}{\large\bfseries\color{accent}}{\thesubsection}{0.6em}{}
\titlespacing*{\section}{0pt}{16pt}{6pt}
\titlespacing*{\subsection}{0pt}{12pt}{4pt}

\pagestyle{fancy}
\fancyhf{}
\fancyhead[L]{\small Model Development Plan}
\fancyhead[R]{\small Alqueva PSP--PV--BESS Trading Optimizer}
\fancyfoot[C]{\small\thepage}
\renewcommand{\headrulewidth}{0.3pt}

\newcommand{\notebox}[2]{\par\smallskip\noindent\colorbox{soft}{\parbox{\dimexpr\linewidth-2\fboxsep}{\textbf{#1} #2}}\par}

\begin{document}

{\LARGE\bfseries\color{accent} Model development, tuning and selection}\\[4pt]
{\large A plan for building the training section the way a data-science team in a company would}\\[2pt]
{\small Written 2026-10-11. Reference document; nothing in Sections 4--6 has been built yet.}

\medskip
\hrule

\section{Why this document}

Two questions came up while preparing for interviews, and one decision followed from them.

\begin{enumerate}
\item \textbf{How can algorithms that have different hyperparameters be compared} (a decision tree against LightGBM, for
      example), and how is one of them chosen to forecast a new day?
\item \textbf{How should hyperparameter tuning be shown} so that it is a strong, defensible part of the project and not just
      a grid search that anyone can run?
\item \textbf{Decision:} build this as a proper section of the project, with the old fixed-settings selection replaced once
      the new one is shown to work.
\end{enumerate}

\section{How different algorithms are compared}

\subsection{The principle}

The hyperparameters are never compared. Each algorithm is a different function with its own knobs (depth and leaf size for
a tree; learning rate, number of leaves and sampling for LightGBM). What is compared is \textbf{how wrong the forecasts
were on days the model had not seen}. A ``model'' in this comparison is the algorithm together with its tuning procedure.

\subsection{The procedure}

\begin{enumerate}
\item \textbf{Tune each algorithm on its own}, using only its training window. An inner validation block (the latest days
      of the training window, separated by a gap) picks the settings. All algorithms get the \textbf{same search effort}.
\item \textbf{Score all tuned algorithms the same way}: rolling-forward folds, same delivery days, same data, same metrics
      (sum of squared errors, MAE and RMSE in price units). Never a random split, which lets the future leak into the past.
\item \textbf{Test whether a difference is real.} The lowest MAE can be luck. The Diebold--Mariano test on daily errors
      says whether one model is better or whether the two cannot be told apart.
\item \textbf{Use a written selection rule.} If the leading models are not significantly different, prefer an average of
      them or the simpler model.
\item \textbf{Score the chosen model once on a locked test period} that played no part in tuning or selection. That is
      the honest accuracy to quote.
\item \textbf{For a new day, use the selected algorithm refitted on the latest data}, and repeat the selection on a
      schedule, because the best algorithm changes with the market.
\end{enumerate}

\subsection{``Which algorithm is best for one particular day?''}

One day cannot answer this: a single day is mostly luck. What can be shown is how often each model wins across many
days, hours and market regimes, and which model has been best recently. The error-weighted combination in the framework
already uses the second idea (weights from the last 28 days of errors, lagged by two days).

\subsection{Weakness to keep in mind}

Equal \emph{number} of trials is not equal \emph{coverage}: LightGBM has eight settings to search and a decision tree
two, so twelve draws cover the large space more thinly. The plan raises the budget and uses a smarter search for this
reason.

\section{Where the project stands}

\subsection{What the old pipeline does}

The existing price forecasters (day-ahead, IDA1--3, XBID, aFRR, mFRR) share \texttt{ml\_\allowbreak train\_\allowbreak val\_\allowbreak test\_\allowbreak common.py}. Each
algorithm (XGBoost, LightGBM, random forest) is fitted with \textbf{fixed, hand-set} parameters, scored on
walk-forward folds, and the lowest mean MAE is chosen. A Wilcoxon test on per-fold errors is available.

\subsection{What was added in Phases 6, 7 and 9}

\begin{itemize}
\item \textbf{Phase 6:} one evaluation contract, rolling-origin backtest, Diebold--Mariano test with HAC variance, locked
      test period, leakage check. \textbf{Phase 7:} point-in-time feature catalogue and store.
\item \textbf{Phase 9:} nested, time-aware tuning (12 random draws, defaults always included, re-tuned every 180 days),
      forecast combinations, and a comparison of window length and refit frequency.
\end{itemize}

\subsection{Phase 9 results on the day-ahead price (development window, 1\,015 days)}

\begin{center}
\begin{tabular}{@{}lrrrr@{}}
\toprule
Model & MAE untuned & MAE tuned & Tuned vs untuned & DM $p$ \\
\midrule
Decision tree     & 20.28 & 19.82 & better & 0.012 \\
Random forest     & 18.11 & 18.06 & no clear difference & 0.46 \\
Gradient boosting & 17.96 & 17.77 & no clear difference & 0.17 \\
XGBoost           & 18.24 & 17.72 & better & 0.0003 \\
LightGBM          & 18.29 & \textbf{17.43} & better & $<0.0001$ \\
\bottomrule
\end{tabular}
\end{center}

\begin{itemize}
\item \textbf{Tuning changes the ranking.} Untuned, gradient boosting led; tuned, LightGBM leads. This is the reason a
      comparison of fixed-setting models is not trustworthy.
\item \textbf{Combinations} (mean, median, trimmed mean, inverse error, top three) are within 0.1 MAE of the best single
      tuned model and cannot be told apart from it ($p$ from 0.22 to 0.99). They are insurance, not an improvement.
\item \textbf{A longer calibration window helped as much as tuning.} LightGBM at defaults: 730 days 18.29, expanding window
      17.56 (DM $p<0.001$). Refit every 7, 30 or 90 days made no detectable difference. Whether tuning and the longer
      window add up is \textbf{not tested}.
\item The tuned run took 7.7 hours while sharing the machine with two other jobs, and 38 minutes when repeated alone; the
      repeat reproduced every score exactly. All 30 tuning rounds beat the defaults on the inner validation.
\end{itemize}

\notebox{Limits.}{One target and the development window only; the locked test period is untouched. Twelve random draws is a
light search. The significance tests are not adjusted for the number of comparisons.}

\section{Plan}

The code goes in \texttt{common\_layer/model\_development/}; a showcase folder \texttt{model\_development/} holds the README
that tells the story, the tuning logs, the figures and one script that reproduces everything. Tests in
\texttt{tests/test\_model\_development.py}. Tree models only (decision tree, random forest, gradient boosting, XGBoost,
LightGBM); metrics SSE, MAE and RMSE, as already decided.

\begin{longtable}{@{}>{\raggedright\arraybackslash}p{0.4cm}>{\raggedright\arraybackslash}p{3.0cm}>{\raggedright\arraybackslash}p{7.0cm}>{\raggedright\arraybackslash}p{4.4cm}@{}}
\toprule
\# & Stage & What is built & Done when \\
\midrule
\endhead
1 & Experiment tracking & A local run registry (SQLite): every run stores model, feature set, window, settings, scores,
data digest and code version. & Every number in the documents traces to a run. \\
2 & Fair, stronger tuning & Same trial budget for every algorithm (for example 40). Smarter search (Optuna). Several
rolling inner validation blocks instead of one 90-day block. Parallel trials. & Tuned never worse than default on
validation; no leakage (existing tests plus new ones). \\
3 & Tuning evidence & Log every trial. Figures: error against each setting, which settings matter, stability of choices
across folds, default against tuned, overfitting gap (inner against outer error), time cost. & The tuning folder reads
without running anything. \\
4 & Algorithm comparison & Five tuned algorithms on the same days. Pairwise Diebold--Mariano matrix with Holm correction.
Win rate by hour, month and regime. & The ranking states which differences are real. \\
5 & Selection policy & A written rule in the config: lowest tuned MAE on the development window; if the top models are not
significantly different, prefer an average or the simpler model. & The rule is applied by code, not by hand. \\
6 & Locked test & The champion is scored once on the locked period; the digest is locked. & One honest number per
target. \\
7 & Champion registry & The fitted model, settings, feature-set version and digests are stored. A challenger replaces the
champion only if it beats it significantly over the last $N$ days. & Live code loads the champion from the registry. \\
8 & Monitoring and retraining & Rolling error tracking, a drift flag, re-tuning every 180 days, refit every 30 days. & A
report shows when the champion degrades. \\
9 & Roll-out & Repeat for DA, IDA1, IDA2, IDA3, XBID, aFRR and mFRR. & One table of champions with locked-test
scores. \\
10 & Replace the old selection & Switch the dashboard and the phase forecasters from fixed-setting, lowest-MAE selection to
the registry. & Existing pipeline tests still pass. \\
\bottomrule
\end{longtable}

\subsection{Order and cost}

Stages 1--6 on the day-ahead price first, then 7--8, then 9, and 10 last, because stage 10 changes the live pipeline.
Compute is the main cost: a tuned run for one target at 12 trials took 38 minutes alone (7.7 hours when three jobs shared
the machine). A larger budget needs parallel trials and re-tuning only every 180 days; the time will be measured before it
is promised.

\section{What happens to the old code}

Nothing is deleted before the new section has been shown to work. The old code runs the live pipeline, and it is the
baseline that the new framework has to beat.

\begin{itemize}
\item \textbf{Replaced at stage 10:} the fixed-settings fitting functions and lowest-MAE selection in
      \texttt{ml\_\allowbreak train\_\allowbreak val\_\allowbreak test\_\allowbreak common.py}; the per-forecaster \texttt{*\_train\_val\_test.py} scripts; and
      \texttt{tests/\allowbreak test\_\allowbreak model\_\allowbreak selection\_\allowbreak metrics.py}, which is rewritten for the new selection.
\item \textbf{Kept:} the forecaster modules (they load a model and predict, but from the registry); the PV and
      reservoir-inflow forecasters (not price targets, outside this programme unless requested); the walk-forward utilities
      until each use has been checked.
\item \textbf{Safe procedure:} build alongside the old code; compare old against new on the same days and record it;
      switch one forecaster at a time, running its tests after each; delete the old code in a separate commit so it can be
      restored from history.
\end{itemize}

\section{Decisions needed}

\begin{enumerate}
\item \textbf{Optuna} as a dependency, or stay with random search? (Recommended: Optuna.)
\item \textbf{Registry:} a local SQLite registry, or MLflow? (Recommended: local first, with a clean path to MLflow.)
\item \textbf{First target:} day-ahead price, or another? (Recommended: day-ahead.)
\item \textbf{Stage 10:} agreement to change the existing phase forecasters at the end.
\end{enumerate}

\section{Short answers for an interview}

\textbf{``How do you compare algorithms with different hyperparameters?''} I do not compare the hyperparameters. I tune each
algorithm on its own with the same search effort, using only training data and a time-ordered validation block, and then I
compare the tuned models by out-of-sample error on the same days. I test the difference with Diebold--Mariano, and if two
models cannot be separated I average them or take the simpler one.

\textbf{``How do you pick the model for a new day?''} The selected algorithm, refitted on the latest data, with a locked
test period that was never used for selection, and a scheduled re-check, because the best algorithm changes with the
market.

\textbf{``Did tuning help?''} For LightGBM and XGBoost, yes (4.7\% and 2.8\% lower MAE, significant). For the forest and
gradient boosting, not significantly. A longer training window helped as much as tuning, and an average of models was as
good as the best single model.

\end{document}
